\documentclass[a4paper]{article}
\usepackage{amsfonts}
\usepackage{amsmath}
\usepackage{amssymb}
\usepackage{graphicx}     

\newcommand{\dint}{\displaystyle\int}
\newcommand{\bdint}{\displaystyle\int\limits_{-\infty}^{+\infty}}

\begin{document}
  
\noindent \large Auteur: Thijs Van Schel \\
\noindent \large Studierichting: Informatica\\
\noindent \large Oefeningenreeks 5B nr. 2.2\\

\medskip

\normalsize

$\bdint \dfrac{dx}{x^2+2x+2}, \exists a \in ]-\infty, +\infty[$\\

$ = \bdint \dfrac{dx}{(x+1)^2 + 1}$\\

$ = \dint\limits_{-\infty}^a \dfrac{dx}{(x+1)^2 + 1} + \dint\limits_a^{+\infty} \dfrac{dx}{(x+1)^2 + 1}$\\

$ = \lim\limits_{b \rightarrow -\infty} \dint\limits_b^a \dfrac{dx}{(x+1)^2 + 1} + \lim\limits_{c \rightarrow +\infty} \dint\limits_a^c \dfrac{dx}{(x+1)^2 + 1}$\\

$ = \lim\limits_{b \rightarrow -\infty}[\operatorname{Bgtan}(x+1)]^a_b + \lim\limits_{c \rightarrow +\infty}[\operatorname{Bgtan}(x+1)]^c_a$\\

$ = \operatorname{Bgtan}(a+1) + \dfrac{\pi}{2} - \operatorname{Bgtan}(a+1) + \dfrac{\pi}{2}$\\

$ = \dfrac{2\pi}{2} = \pi$\\

\begin{thebibliography}{999}
%\bibitem{a} Referentie
\end{thebibliography}
\end{document} 
